% Generated by thesis/ch7/exp2/signal_ticks.py -- do not edit by hand.
\begin{table}[p]
  \centering
  \footnotesize
  \setlength{\tabcolsep}{3pt}
  \begin{tabularx}{\textwidth}{@{}>{\raggedright\arraybackslash}X*{10}{c}c@{}}
    \toprule
    Signal & s0 & s1 & s2 & s3 & s4 & s5 & s6 & s7 & s8 & s9 & $n$ \\
    \multicolumn{12}{@{}l}{\textit{Linear core (slopes and one-hot offsets), strongest in the data first}} \\
    \textbf{Marital status and household role}: low income likelier outside marriage: husbands and wives sit well below the unmarried & \checkmark & \checkmark & \checkmark & \checkmark & \checkmark & \checkmark & \checkmark & \checkmark & \checkmark & \checkmark & 10 \\
    \textbf{Occupation and work class}: executive, managerial and professional work sit below; service and handling work and missing job records above & \checkmark & \checkmark & \checkmark & \checkmark & \checkmark & \checkmark & \checkmark & \checkmark & \checkmark & \checkmark & 10 \\
    \textbf{Education (years)}: each extra year of schooling lowers the chance of a low income & \checkmark & \checkmark & \checkmark & \checkmark & \checkmark & \checkmark & \checkmark & \checkmark & \checkmark & \checkmark & 10 \\
    \textbf{Capital gain (slope)}: capital gains lower the chance of a low income & \checkmark & \checkmark & \checkmark & \checkmark & \checkmark & \checkmark & \checkmark & \checkmark & \checkmark & \checkmark & 10 \\
    \textbf{Age (slope)}: older people are less likely to have a low income & \checkmark & \checkmark & \checkmark & \checkmark & \checkmark & \checkmark & \checkmark & \checkmark & \checkmark & \checkmark & 10 \\
    \textbf{Hours per week (slope)}: longer working weeks lower the chance of a low income & \checkmark & \checkmark & \checkmark & \checkmark & \checkmark & \checkmark & \checkmark & \checkmark & \checkmark & \checkmark & 10 \\
    \textbf{Capital loss (slope)}: recorded capital losses lower the chance of a low income & \checkmark & \checkmark & \checkmark & \checkmark & \checkmark & \checkmark & \checkmark & \checkmark & \checkmark & \checkmark & 10 \\
    \textbf{Sex}: women are more likely to have a low income than otherwise similar men & \checkmark & \checkmark & \checkmark & \checkmark & \checkmark & \checkmark & \checkmark & \checkmark & \checkmark & \checkmark & 10 \\
    \textbf{Race}: small offsets by recorded race & \checkmark\textsuperscript{a} & \checkmark\textsuperscript{a} & \checkmark & \checkmark\textsuperscript{a} & \checkmark\textsuperscript{a} & \checkmark\textsuperscript{a} & \checkmark\textsuperscript{a} & \checkmark\textsuperscript{a} & \checkmark\textsuperscript{a} & \checkmark\textsuperscript{a} & 10 \\
    \textbf{Education levels}: offsets for single education levels beyond the years-of-schooling line & \checkmark\textsuperscript{a} & \checkmark & \checkmark & \checkmark & \checkmark & \checkmark & \checkmark & \checkmark & \checkmark & \checkmark & 10 \\
    \textbf{Country of birth}: offsets for single countries of birth, mostly small groups & \checkmark\textsuperscript{a} & \checkmark & \checkmark & \checkmark & \checkmark & \checkmark & \checkmark\textsuperscript{a} & \checkmark & \checkmark\textsuperscript{a} & \checkmark & 10 \\
    \midrule
    \multicolumn{12}{@{}l}{\textit{Non-linear: bends and interactions, strongest in the data first}} \\
    \textbf{Capital-gain threshold}: past a few thousand dollars of capital gain the chance of a low income drops steeply \emph{(usually split across several units)} & \checkmark\textsuperscript{a} & \checkmark & \checkmark & \checkmark &  & \checkmark & \checkmark\textsuperscript{b} & \checkmark\textsuperscript{c} & \checkmark & \checkmark & 9 \\
    \textbf{Age bend}: the fall with age is steep for the young and flattens or turns back later in life \emph{(usually a single gate that also reads another input)} &  &  &  & \checkmark\textsuperscript{b} & \checkmark & \checkmark\textsuperscript{d} & \checkmark\textsuperscript{b} & \checkmark & \checkmark &  & 6 \\
    \textbf{Hours bend}: short weeks raise the chance of a low income more than long weeks lower it \emph{(usually a single hinge)} &  &  &  &  &  & \checkmark &  &  &  & \checkmark & 2 \\
    \midrule
    \multicolumn{12}{@{}l}{\textit{Relevant in the data but in no network}} \\
    \textbf{Capital-loss bend}: bends in the capital-loss effect across loss amounts &  &  &  &  &  &  &  &  &  &  & 0 \\
    \textbf{Census weight (slope)}: the survey sampling weight, not a personal attribute &  &  &  &  &  &  &  &  &  &  & 0 \\
    \textbf{Age $\times$ marriage}: the age effect differs inside and outside marriage &  &  &  &  &  &  &  &  &  &  & 0 \\
    \textbf{Census weight bend}: bends in the survey sampling weight &  &  &  &  &  &  &  &  &  &  & 0 \\
    \textbf{Age $\times$ hours}: the hours effect changes with age &  &  &  &  &  &  &  &  &  &  & 0 \\
    \textbf{Capital gain $\times$ marriage}: the capital-gain threshold acts differently for the married &  &  &  &  &  &  &  &  &  &  & 0 \\
    \textbf{Sex $\times$ marriage}: the gap between women and men differs inside and outside marriage &  &  &  &  &  &  &  &  &  &  & 0 \\
    \bottomrule
  \end{tabularx}
  \caption{Adult: signals (rows) found in each seed's network (columns). \checkmark: the signal appears in that seed. A lettered note marks a seed that realises the signal differently from most: \textsuperscript{a}\,present but weak (0.03 to 0.10 logit); \textsuperscript{b}\,a single hinge; \textsuperscript{c}\,a single gate that also reads another input; \textsuperscript{d}\,split across several units. The positive class is an income of at most \$50K.}
  \label{tab:exp2-signalticks-adult}
\end{table}
